\section{Crítica ao Código Original}

\label{sec:critica}

O \texttt{Main.java} fornecido implementa corretamente o funcionamento do jogo,
porém sua estrutur segue um modelo inteiramente procedural, apesar de ter sido
desenvolvido em Java. Conceitos fundamentais da programação orientada a objetos,
como encapsulamento, abstração, herança, polimorfismo e o uso de coleções,
praticamente não são explorados. Em vez disso, toda a lógica da aplicação é
concentrada em um único método, utilizando variáveis primitivas e vetores para
representar as entidades do jogo. Em suma, pode ser comparado a um código de alguém que sabe o básico
de Linguagem C e decidiu se aventurar no Java.

Embora essa abordagem seja suficiente para um projeto pequeno, ela dificulta a
organização do código, reduz seu potencial de reutilização e torna sua evolução
mais trabalhosa.

\subsection{Ausência de modelagem orientada a objetos}

Toda a lógica do jogo está concentrada na classe \texttt{Main},
principalmente dentro do método \texttt{main()}, auxiliado apenas por alguns
métodos utilitários estáticos. Não existe uma modelagem que represente os
conceitos do domínio do jogo, como jogador, inimigos, projéteis ou explosões,
mesmo sendo entidades que possuem características e comportamentos próprios.

Como consequência, o estado de cada entidade é representado por variáveis
primitivas independentes, manipuladas diretamente ao longo do programa, como no
caso do jogador:

\begin{lstlisting}[caption={Estado do jogador: variáveis primitivas sem uma classe que represente a entidade (Main.java)}]
int player_state = ACTIVE;
double player_X = GameLib.WIDTH / 2;
double player_Y = GameLib.HEIGHT * 0.90;
double player_VX = 0.25;
double player_VY = 0.25;
double player_radius = 12.0;
double player_explosion_start = 0;
double player_explosion_end = 0;
long player_nextShot = currentTime;
\end{lstlisting}

Essa representação impede que os dados e o comportamento relacionados ao
jogador sejam reunidos em uma única estrutura. Como não há encapsulamento,
qualquer trecho do \texttt{main()} pode acessar ou modificar diretamente
variáveis como \texttt{player\_X}, \texttt{player\_Y} ou
\texttt{player\_state}, sem qualquer mecanismo de controle ou validação.

Além disso, concentrar toda a lógica da aplicação em um único método faz com
que responsabilidades bastante diferentes, como processamento de entrada,
movimentação, geração de entidades, detecção de colisões e desenho,
permaneçam fortemente interligados. Essa organização dificulta a manutenção do
código e aumenta o impacto de qualquer simples alteração realizada no sistema.

\subsection{Representação das entidades por vetores paralelos}

A ausência de objetos torna-se ainda mais evidente nas entidades que possuem
múltiplas instâncias, como inimigos e projéteis. Em vez de cada inimigo ser
representado por um objeto contendo seus próprios atributos, o programa utiliza
diversos vetores paralelos, nos quais cada posição corresponde implicitamente à
mesma entidade.

O inimigo do tipo 1, por exemplo, é representado pelos seguintes vetores:

\begin{lstlisting}[caption={Inimigo tipo 1 representado por vetores paralelos (Main.java)}]
int    [] enemy1_states          = new int[10];
double [] enemy1_X               = new double[10];
double [] enemy1_Y               = new double[10];
double [] enemy1_V               = new double[10];
double [] enemy1_angle           = new double[10];
double [] enemy1_RV              = new double[10];
double [] enemy1_explosion_start = new double[10];
double [] enemy1_explosion_end   = new double[10];
long   [] enemy1_nextShoot       = new long[10];
\end{lstlisting}

O mesmo padrão é repetido para
os demais tipos de inimigo. Dessa forma, atributos que pertencem
conceitualmente ao mesmo objeto ficam espalhados em diferentes estruturas,
dependendo do índice do vetor para manter sua associação.

Essa escolha também impede o reaproveitamento de código entre entidades que
compartilham características semelhantes. Conceitos comuns, como posição,
velocidade, estado e remoção ao sair da tela, acabam sendo implementados
repetidamente para cada conunto de vetores, em vez de serem centralizados em
uma classe base.

Outro efeito dessa representação é a necessidade de gerenciar manualmente as
posições disponíveis dentro dos vetores. Sempre que uma nova entidade precisa
ser criada, o programa realiza uma busca linear para localizar um índice livre,
como mostrado a seguir:

\begin{lstlisting}[caption={Busca manual por um índice livre no vetor (Main.java)}]
public static int findFreeIndex(int [] stateArray){
    int i;
    for(i = 0; i < stateArray.length; i++){
        if(stateArray[i] == INACTIVE) break;
    }
    return i;
}
\end{lstlisting}

Na prática, esse método reproduz manualmente uma funcionalidade que já é
oferecida pelas coleções da biblioteca padrão do Java. Usando estruturas como
\texttt{List}, por exemplo, acabos por eliminar a necessidade de procurar
posições livres e de manter índices sincronizados entre diversos vetores.

Além disso, o uso de vetores de tamanho fixo impõe uma limitação artificial ao
número de entidades existentes simultaneamente no jogo. Quando todos os índices
estão ocupados, novas entidades simplesmente deixam de ser criadas.

\subsection{Consequências da arquitetura adotada}

A representação das entidades por meio de variáveis primitivas e vetores
paralelos afeta diretamente toda a implementação do jogo. Como não existe uma
abstração comum entre os diferentes tipos de entidades, operações semelhantes
precisam ser reimplementadas para cada caso específico, resultando em grande
quantidade de código duplicado (o que nunca é positivo).

Esse problema aparece na detecção de colisões. O programa
possui cinco laços distintos para verificar colisões entre diferentes pares de
entidades, como jogador e projéteis inimigos, jogador e inimigos, ou projéteis
do jogador contra inimigos. Embora trabalhem sobre vetores diferentes, todos
executam exatamente a mesma sequência de operações (calcular a distância entre
os objetos e compará-la com a soma de seus raios).

\begin{lstlisting}[caption={Dois dos cinco laços de colisão presentes no código original (Main.java)}]
for(int i = 0; i < e_projectile_states.length; i++){
    double dx = e_projectile_X[i] - player_X;
    double dy = e_projectile_Y[i] - player_Y;
    double dist = Math.sqrt(dx * dx + dy * dy);
    if(dist < (player_radius + e_projectile_radius) * 0.8){
        player_state = EXPLODING;
        player_explosion_start = currentTime;
        player_explosion_end = currentTime + 2000;
    }
}

for(int i = 0; i < enemy1_states.length; i++){
    double dx = enemy1_X[i] - player_X;
    double dy = enemy1_Y[i] - player_Y;
    double dist = Math.sqrt(dx * dx + dy * dy);
    if(dist < (player_radius + enemy1_radius) * 0.8){
        player_state = EXPLODING;
        player_explosion_start = currentTime;
        player_explosion_end = currentTime + 2000;
    }
}
\end{lstlisting}

A única diferença entre esses trechos está na origem dos dados utilizados. A
lógica da colisão permanece exatamente a mesma, sendo apenas adaptada para cada
conjunto de vetores adequedo. Uma representação baseada em objetos permitiria aplicar o
mesmo algoritmo de colisão a qualquer entidade que fornecesse informações como
posição e raio, eliminando essa repetição.

O mesmo padrão também aparece na atualização das entidades e no desenho da
cena. A movimentação dos projéteis do jogador e dos projéteis inimigos, por
exemplo, é praticamente idêntica, alterando apenas a direção em que deixam a
tela e os vetores utilizados.

\begin{lstlisting}[caption={Atualização dos projéteis do jogador e dos projéteis inimigos (Main.java)}]
for(int i = 0; i < projectile_states.length; i++){
    if(projectile_states[i] == ACTIVE){
        if(projectile_Y[i] < 0) {
            projectile_states[i] = INACTIVE;
        } else {
            projectile_X[i] += projectile_VX[i] * delta;
            projectile_Y[i] += projectile_VY[i] * delta;
        }
    }
}

for(int i = 0; i < e_projectile_states.length; i++){
    if(e_projectile_states[i] == ACTIVE){
        if(e_projectile_Y[i] > GameLib.HEIGHT) {
            e_projectile_states[i] = INACTIVE;
        } else {
            e_projectile_X[i] += e_projectile_VX[i] * delta;
            e_projectile_Y[i] += e_projectile_VY[i] * delta;
        }
    }
}
\end{lstlisting}

Situações semelhantes ocorrem em diversos outros pontos do programa, como no
desenho das entidades e na atualização dos diferentes tipos de inimigo. Sempre
que um novo tipo de objeto é introduzido, torna-se necessário escrever novos
laços muito parecidos com os já existentes, alterando apenas nomes de vetores e
algumas constantes.

\subsection{Impacto na manutenção e evolução do sistema}

Além dos problemas de organização, a estrutura do código limita a evolução do
jogo. Como toda a lógica está concentrada em um único método e cada tipo de
entidade possui sua própria implementação, qualquer funcionalidade nova exige
alterações em diversos pontos diferentes do programa.

Essa dificuldade torna-se ainda mais evidente no gerenciamento das entidades.
Todos os vetores possuem tamanho fixo, definido manualmente durante sua criação.
Quando não há mais posições disponíveis, novas entidades simplesmente deixam de
ser criadas.

\begin{lstlisting}[caption={Disparo descartado quando todos os índices do vetor estão ocupados (Main.java)}]
int free = findFreeIndex(projectile_states);
if(free < projectile_states.length){
    projectile_X[free] = player_X;
    projectile_Y[free] = player_Y - 2 * player_radius;
    projectile_VX[free] = 0.0;
    projectile_VY[free] = -1.0;
    projectile_states[free] = ACTIVE;
    player_nextShot = currentTime + 100;
}
// se free == projectile_states.length, o tiro simplesmente não acontece
\end{lstlisting}

Essa limitação obriga o programa a manter capacidades máximas definidas
arbitrariamente e a controlar manualmente quais posições dos vetores estão
livres ou ocupadas. Em aplicações maiores, esse tipo de gerenciamento aumenta a
complexidade do código e dificulta a expansão do sistema.

Por fim, praticamente todas as responsabilidades do jogo permanecem concentradas dentro
do método \texttt{main()}. A ausência de uma divisão clara entre essas
responsabilidades reduz a modularidade da aplicação e faz com que alterações em
uma funcionalidade tenham maior probabilidade de impactar outras partes do
sistema.


